\documentclass[11pt]{article}
\usepackage[a4paper,margin=25mm,headheight=14pt]{geometry}
\usepackage[T1]{fontenc}
\usepackage{lmodern,microtype}
\usepackage{amsmath,amssymb,amsthm,mathtools}
\usepackage{booktabs,longtable,array}
\usepackage{xcolor,enumitem,fancyhdr}
\usepackage[colorlinks=true,linkcolor=blue!50!black,urlcolor=blue!50!black,citecolor=blue!50!black]{hyperref}
\usepackage{listings}
\hypersetup{pdftitle={Weak divisibility and shared sign certificates for Diophantine consistency},pdfauthor={Research prepared for Vladimir Reshetnikov with AI assistance},pdfsubject={Explicit integer polynomials for Con(ZFC); eleven variables and direct quartics}}
\newcommand{\Z}{\mathbb Z}
\newcommand{\N}{\mathbb N}
\newcommand{\Q}{\mathbb Q}
\newcommand{\Con}{\operatorname{Con}}
\newcommand{\ZFC}{\mathsf{ZFC}}
\newcommand{\sgn}{\operatorname{sgn}}
\newcommand{\vp}{\operatorname{v}}
\newcommand{\Norm}{\mathcal N_3}
\newcommand{\src}{\texttt{source34}}
\newtheorem{theorem}{Theorem}[section]
\newtheorem{lemma}[theorem]{Lemma}
\newtheorem{proposition}[theorem]{Proposition}
\newtheorem{corollary}[theorem]{Corollary}
\theoremstyle{definition}
\newtheorem{definition}[theorem]{Definition}
\theoremstyle{remark}
\newtheorem{remark}[theorem]{Remark}
\setlist{itemsep=3pt,topsep=6pt}
\setlength{\emergencystretch}{3em}
\setcounter{tocdepth}{1}
\pagestyle{fancy}
\fancyhf{}
\fancyhead[L]{\small Diophantine consistency}
\fancyhead[R]{\small Weak carriers and shared signs}
\fancyfoot[C]{\thepage}
\lstset{basicstyle=\ttfamily\small,breaklines=true,columns=fullflexible,frame=single,rulecolor=\color{black!20}}

\title{Weak divisibility and shared sign certificates\\for Diophantine consistency}
\author{A research contribution prepared for Vladimir Reshetnikov's ProveIt project\\\small With AI assistance; ordinary proofs and reproducible arithmetic artifacts}
\date{24 October 2026}

\begin{document}
\maketitle
\begin{abstract}
We give new explicit integer polynomials whose lack of integer solutions
is, using the repository's stated ordinary calibration theorems, provably
equivalent in ZFC to $\Con(\ZFC)$. The principal small-variable equation
has eleven unrestricted integer unknowns, total degree $416$, and $129$
live binary arithmetic operations. A companion has degree $464$ and
$123$ operations. The closest inspected incoming eleven-integer equation
has degree $688$ and $126$ operations. We also give a direct quartic in
$478$ unrestricted integers, built from the numerical proof calculus
C222 without K215, Jones's reduction, or a large fixed coefficient tuple.
Its $305$ quadratic residuals expand to $5{,}484$ monomials, with maximum
absolute coefficient $511{,}326$.

The first improvement combines a weak fifth-power congruence, a mixed
modulus carrier, coprime fusion, and interpolation with a fixed strong
Pell witness. Three-adic root separation removes an unnecessary even
multiplier and permits a cheaper center. The direct improvement proves
some source values positive earlier, allowing pairs of sign conditions
to share one Diophantine witness. A joint guard also handles the signs
and nonvanishing needed by the shared Pell frame. Full coefficient files,
literal polynomial graphs, deterministic builders, and independent
finite checks accompany the ordinary proofs. The work does not claim
global minimality, a proof of consistency, or a new end-to-end
proof-assistant derivation.
\end{abstract}

\tableofcontents
\newpage

\section{Results, comparison, and logical scope}
\label{sec:results}

The problem is to specify a fixed polynomial $P\in\Z[x_1,\ldots,x_m]$
and prove
\begin{equation}\label{eq:goal}
 \ZFC\vdash
 \left(\forall\mathbf x\in\Z^m\; P(\mathbf x)\ne0\right)
 \longleftrightarrow\Con(\ZFC).
\end{equation}
All variables in this paper are unrestricted integers unless explicitly
stated otherwise. A polynomial abbreviation is a shared subexpression,
not an extra existentially quantified coordinate. Fixed numeral ports
are actual supplied integers, not free parameters.

There is no single agreed ordering of simplicity: variable count,
degree, arithmetic length, coefficient height, and the complexity of
the calibration can favor different equations. We give measured
improvements on two existing routes. We do not give a global lower
bound on possible constructions. In particular, ten-integer
constructions already occur in the repository, and a recent general
ten-unknown reduction is available in the literature\cite{sun2026}.
An eleven-variable degree improvement is not a new minimum-variable
theorem.

\begin{table}[ht]
\centering
\begin{tabular}{lrrr}
\toprule
Eleven-integer construction & Variables & Degree & Operations\\
\midrule
Incoming coprime fusion, centered & 11 & 688 & 126\\
Incoming coprime fusion, direct & 11 & 728 & 122\\
Incoming all-radicand fusion & 11 & 1632 & 114\\
New cubic chart, direct & 11 & 488 & 120\\
New cubic chart, centered & 11 & 464 & 123\\
New mixed chart, direct & 11 & 456 & 126\\
\textbf{New mixed chart, centered} & \textbf{11} & \textbf{416} & \textbf{129}\\
\bottomrule
\end{tabular}
\caption{All four new equations use the same complete six-port tuple.
Operations count every live binary $+,-,\times$, including scalar
products and squares; leaves and reuse are free.}
\label{tab:norm}
\end{table}

Degree $416$ is $39.53\%$ below $688$, at three additional operations.
The $464/123$ companion improves both degree and arithmetic length over
the $688/126$ predecessor. The $456/126$ companion keeps that predecessor's
operation count and lowers its degree by $33.72\%$.

\begin{table}[ht]
\centering
\begin{tabular}{lrrrr}
\toprule
Direct integer quartic & Variables & Rows & Terms & Height\\
\midrule
Matched incoming derived-label form & 517 & 310 & 6774 & 3904267\\
\textbf{New shared-sign form} & \textbf{478} & \textbf{305} & \textbf{5484} & \textbf{511326}\\
\bottomrule
\end{tabular}
\caption{Rows are quadratic residuals, terms are collected monomials
of their sum of squares, and height is the largest absolute coefficient.}
\label{tab:quartic}
\end{table}

The direct result improves all four displayed measures of its matched
predecessor. Other incoming quartics trade more variables for smaller
coefficients; this table is not a claim of dominance over every known
quartic on every measure.

\paragraph{Source snapshot.}
The comparison snapshot is ProveIt commit \texttt{45b399ec51df};
the complete identifier is recorded in \texttt{SOURCE\_SNAPSHOT.json}.
We inspected the
current research comparison, the numbered collection, and all eight
incoming archives at that revision. The immediate ingredients are the
incoming \emph{Coprime Projection}, \emph{Powered Saturation and
Factorial Carriers}, \emph{Carriers and Even Indices}, and
\emph{Quadratic Frames} packages\cite{repo,incoming}. The first three
already contain the individual weak-congruence, mixed-carrier, and
fusion ideas used here. Our claim concerns their new compatible
composition, its smaller literal equations, and the direct sign
reductions proved below. We make no literature-wide priority claim for
the component tricks.

\paragraph{What is proved and what is inherited.}
There are three distinct levels of evidence.
\begin{enumerate}
\item The new arithmetic reductions are proved by ordinary mathematics,
      including all zero branches, negative signs, and the order of
      existential witness selection.
\item The literal files are checked by exact finite algebra, authenticated
      fixed numerals, exact expansion or degree certificates, and
      deterministic reconstruction. These checks do not prove the
      unbounded mathematical theorems.
\item The consistency interpretation uses the named ordinary source
      calibration. The small-variable route uses the native compiler
      and the K215 endpoint. The direct route uses C222 refutability,
      its ZF calibration, and the ZF--ZFC refutation equivalence.
      Inherited Lean endpoint texts identify these interfaces, but no
      full Lean/Rocq dependency closure is rebuilt in this release.
\end{enumerate}
The unconditional new reduction can therefore also be read independently
as a zero-existence equivalence with the cited predecessor arithmetic
source. Composing with that source's proved ordinary calibration gives
\eqref{eq:goal}. The proof uses no assumption that ZFC is consistent.

\section{The complete eleven-integer equation}
\label{sec:equation}

Fix the six integers
\[
 \beta,K_0,M_C,M_F,A_0,B_0
\]
specified in Section~\ref{sec:constants}. In particular
\[
 \beta=2^{1064921}-1,\qquad A_0=2^{1113},\qquad B_0=2^{666104}.
\]
The high mask $M_F$ is already shifted. No coefficient is selected by
waiting for a proof or a computation to halt.

The eleven independent variables are
\begin{equation}\label{eq:variables}
 (j,C,z,h,r,v,w,T,c,n,s)\in\Z^{11}.
\end{equation}
For the principal \emph{mixed chart}, define
\begin{align}
 J&=6j+1,&q&=\beta J+1,&W&=12w+4,\notag\\
 X&=qW,&Y&=qv,&E&=XY,\notag\\
 a&=E+Y,&H&=4a+3,&\Delta&=a^2+H=(a+2)^2-1.\label{eq:front}
\end{align}
Define the packing and budget polynomials
\begin{align}
 F&=(K_0+W)C-z(q-1),\notag\\
 R&=\bigl(q(q-F)-C\bigr)(q^2-1)+(M_C+qM_F)J,\notag\\
 k&=R+1+hE,\qquad \eta=c-Yk,\notag\\
 E_0&=(A_0F-q)^2+(B_0C-q)^2,\notag\\
 B&=k^2(q^2-E_0)-4q^2\eta^2.\label{eq:budget}
\end{align}
The two fused conditions and the strong and auxiliary conditions are
\begin{align}
 \ell&=XY^2,&L&=\ell k,&D_t&=1+L(L+k),\notag\\
 D_b&=3B-4,&D_g&=1+D_bD_ts^2,&P&=D_tD_g,\notag\\
 I&=rc^2,&D_f&=1+\Delta I^2,&Q&=\Delta^2I^2,\notag\\
 \delta&=Q-1,&V&=D_f(c-R+4IT)-c,\notag\\
 C_y&=QV^2-1,&D_a&=\delta C_y.\label{eq:pell}
\end{align}
The common scalar projection and mixed carrier are
\begin{align}
 K&=X(3c-2X)+2,& M&=2q^2,&D&=MH,\notag\\
 C_*&=MK+H(v-q^2),&z_*&=DT+C_*,\notag\\
 A_*&=(z_*^2+3D^2)(z_*^3-19D^3),&S&=srM^5H.\label{eq:carrier}
\end{align}
Finally put
\begin{align}
 Z&=S(n+3QV)-A_*,\notag\\
 u&=S^2D_f,&b&=9S^2D_a,&e&=81S^2P.\label{eq:radicands}
\end{align}
The promised single equation is
\begin{equation}\label{eq:principal}
 \boxed{\quad
 \left[(Z^2-u-b-e)^2-4(ub+ue+be)\right]^2
       -64ubeZ^2=0.
 \quad}
\end{equation}
Every displayed symbol is a polynomial abbreviation in
\eqref{eq:variables}. There are no variable exponents, square roots,
divisions, bounded quantifiers, or hidden witnesses in this definition.
The expression has integer coefficients and is evaluated by the supplied
$129$-operation graph.

\paragraph{Three companion equations.}
Dropping $3QV$ from the definition of $Z$ gives the degree-$456$ direct
mixed chart. It is related to the principal form by an integral
translation of $n$, since $QV$ does not depend on $n$.
For the \emph{cubic chart}, replace $Y=qv$ by
\[
 Y=(2v+1)q^3,
\]
and replace the carrier definitions by
\begin{equation}\label{eq:cubiccarrier}
 D=H,\qquad z_*=HT+K,\qquad
 A_*=(z_*^2+3H^2)(z_*^3-19H^3),\qquad S=srH.
\end{equation}
Keep the other definitions unchanged. With center $3QV$ the degree is
$464$; without it the degree is $488$.

\section{Three elementary mechanisms}
\label{sec:mechanisms}

\subsection{An exact norm with only three roots}
Define the integer polynomial
\begin{equation}\label{eq:norm3}
 \Norm(Z;u,b,e)=
 \left[(Z^2-u-b-e)^2-4(ub+ue+be)\right]^2-64ubeZ^2.
\end{equation}
In a formal algebra with square roots adjoined,
\begin{equation}\label{eq:normproduct}
 \Norm(Z;u,b,e)=
 \prod_{\epsilon\in\{\pm1\}^3}
 \left(Z-\epsilon_1\sqrt u-\epsilon_2\sqrt b-
                   \epsilon_3\sqrt e\right).
\end{equation}
Independent sign invariance removes odd radical powers. The identity
is valid for zero, negative, or dependent radicands as well as for
positive ones. It has $35$ collected monomials, all checked exactly in
the release. A zero of \eqref{eq:principal} therefore makes at least
one of the eight factors vanish over $\mathbb C$.

\begin{lemma}[Three-adic separation]\label{lem:separation}
Let $A,B,C$ be positive integers with $A\equiv B\equiv1\pmod3$.
If a signed sum
\[
 \epsilon_1\sqrt A+3\epsilon_2\sqrt B+9\epsilon_3\sqrt C
\]
is rational, all three roots are integers.
\end{lemma}
\begin{proof}
Write each radicand as a square times a positive squarefree integer.
Distinct squarefree radicals are linearly independent over $\Q$.
The first two squarefree parts are not divisible by $3$, and their
integer coefficients have respective $3$-adic orders $0$ and $1$.
If the third squarefree part is divisible by $3$, its irrational term
cannot join either of them. Otherwise its coefficient has order at
least $2$. Within every nontrivial squareclass, the term of least
$3$-adic order is unique and cannot cancel. No such class can occur.
Rational square roots of integers are integers.
\end{proof}

This separation lets us use $I=rc^2$. The preceding parity separation
used an even multiplier $2src^2$. We need neither its factor $2$ nor
its explicit factor $s$ in the new strong ordinate. Completeness will
choose $I$ divisible by a prescribed multiple of $s$ without making
that factor part of the polynomial.

\subsection{Coprime fusion of the first norm and the guard}
For any integers in the displayed construction,
\begin{equation}\label{eq:coprime}
 \gcd(D_t,D_g)=\gcd(D_t,1+D_bD_ts^2)=1.
\end{equation}
Thus if $D_t,D_g$ are positive and their product $P$ is a square,
each factor is a square. This is the inherited selective fusion idea;
we use it together with the weaker carrier.

There is also a small literal saving in the guard. Since $q,A_0,B_0$
are even, $E_0$ and $B$ are divisible by $4$. Hence
\begin{equation}\label{eq:guardparity}
 D_b=3B-4\equiv0\pmod4,\qquad D_g\equiv1\pmod4.
\end{equation}
The predecessor's additional factor $4$ in $D_g-1$ is unnecessary.
This conclusion uses actual coefficient parity; it is not a statement
for arbitrary replacements of the six fixed ports.

\subsection{A fifth-power congruence is enough}
Write
\begin{equation}\label{eq:pi}
 \pi(t)=(t^2+3)(t^3-19),\qquad
 \Phi_D(z)=D^5\pi(z/D)
 =z^5+3D^2z^3-19D^3z^2-57D^5.
\end{equation}
The familiar intersective polynomial $\pi$ has a root modulo every
positive integer and no rational root\cite{hyde2014}. For completeness,
when $p\ne2,3,19$ either $p\equiv2\pmod3$ and cubing is bijective,
or $p\equiv1\pmod3$ and $-3$ is a quadratic residue. Selected roots
are simple and lift. At $p=19$ the quadratic has a simple root; at
$p=2$ the cubic does. At $p=3$, write $t=1+3u$; the equation
$(1+3u)^3=19$ reduces to
\[
 u+3u^2+3u^3=2,
\]
which has a simple root modulo $3$ and lifts to all powers. CRT
finishes. The quadratic factor has no real zero and $19$ is not a
rational cube, proving the no-rational-root assertion.

\begin{lemma}[Strong and weak parts of one carrier]\label{lem:carrier}
Let $M,H$ be nonzero coprime integers, let $D=MH$, and let
\[
 z=DT+MK+H\omega.
\]
If $M^5H\mid\Phi_D(z)$, then
\[
 M\mid\omega,\qquad H\mid K^5.
\]
\end{lemma}
\begin{proof}
For a prime with $p^e\mid M$ and $v_p(M)=e$, suppose $v_p(z)=b<e$.
The leading term $z^5$ has order $5b$, while each other term in
\eqref{eq:pi} has strictly greater order. Thus $M^5$ cannot divide
$\Phi_D(z)$, a contradiction. Consequently $M\mid z$, and
coprimality gives $M\mid\omega$. Reducing modulo $H$ gives
$H\mid(MK)^5$; cancel the unit $M$ modulo $H$.
\end{proof}

\begin{lemma}[Height resolves weak divisibility]\label{lem:power5}
Suppose $x\ge2^{64}$, $t\ge24$, and
\[
 |H|>x(x-3)^t,\qquad
 |\Pi|<2^{4t+4}x^2.
\]
If $H\mid\Pi^5$, then $\Pi=0$.
\end{lemma}
\begin{proof}
Since $x-3>x/2$,
\[
 |H|>\frac{x^{t+1}}{2^t},\qquad
 |\Pi|^5<2^{20t+20}x^{10}.
\]
The sufficient inequality is
$x^{t-9}>2^{21t+20}$. Its base-two exponent margin is at least
\[
 64(t-9)-(21t+20)=43t-596\ge436.
\]
Thus $|\Pi|^5<|H|$. The only multiple of $H$ of that magnitude is zero.
\end{proof}

The new scale spends the very large actual coefficient bound on this
weaker congruence. It never makes the false inference
$H\mid K^5\Rightarrow H\mid K$ for arbitrary integers.

\section{Soundness of the eleven-integer construction}
\label{sec:soundness}

We prove the mixed chart first. The cubic chart omits the recovery
of its scale, since its scale is already present in the definition.

\subsection{Uniform facts and the zero branches}
The fixed $\beta$ is odd and $\beta\equiv1\pmod3$. Thus $q$ is
even, $q\equiv2\pmod3$, and $|q|\ge\beta-1$. Also
\[
 W\equiv1\pmod3,\quad X\equiv2\pmod3,\quad
 |X|\ge4(\beta-1)>2^{64}.
\]
If $v=0$, then $a=0$, $H=3$, and $\Delta=3$. Otherwise $|a|$
is enormous. In every case $H$ is nonzero and odd, $\Delta>1$, and
\[
 a\equiv H\equiv\Delta\equiv0\pmod3.
\]
Because $Y=qv$, $H\equiv3\pmod q$. As $3\nmid q$,
\begin{equation}\label{eq:MHcoprime}
 \gcd(M,H)=1,\qquad M=2q^2.
\end{equation}
We have $D_f\ge1$ and
$D_t=1+\ell(\ell+1)k^2\ge1$. Also $Q$ is either $0$ or
greater than $1$, so $\delta$ never vanishes.

Suppose \eqref{eq:principal} holds. If $S=0$, all scaled radicands
vanish and the equation becomes $A_*^8=0$. But
\[
 A_*=D^5\pi(T+C_*/D)\ne0
\]
by the no-rational-root property and $D\ne0$. Hence $S\ne0$,
and in particular $s,r\ne0$.

We do \emph{not} yet assume $c\ne0$. If $c=0$, then $I=Q=0$
and $D_a=1$. If $c\ne0$, then $|I|\ge c^2$ and
$D_f>|c|+1$. Writing $b_0=c-R+4IT$, either $b_0=0$ and
$|V|=|c|\ge1$, or $|V|\ge D_f-|c|>1$. Therefore $D_a>0$
in both cases. The roots associated with $D_f$ and $D_a$ are real.

If $D_g<0$, the $P=D_tD_g$ root would be the sole nonzero imaginary
term in a vanishing conjugate of \eqref{eq:normproduct}, impossible.
Equation~\eqref{eq:guardparity} excludes $D_g=0$. Thus $D_g>0$.
Since $B$ is a multiple of $4$, $B\le0$ would imply $D_b\le-4$
and $D_g\le-3$. We conclude
\begin{equation}\label{eq:budgetpositive}
 B\ge4,\quad k\ne0,\quad E_0<q^2,\quad
 \left|\frac{c}{k}-Y\right|<\frac12.
\end{equation}

\subsection{Extract roots before using the restored scale}
Modulo $3$ we have $D_f\equiv D_a\equiv1$. Divide a vanishing
conjugate by $S$. Its radical sum is
\[
 \epsilon_f\sqrt{D_f}+3\epsilon_a\sqrt{D_a}
              +9\epsilon_p\sqrt{D_tD_g}\in\Q.
\]
Lemma~\ref{lem:separation} makes all three roots integral.
Equation~\eqref{eq:coprime} then yields integers $f,\tau,g,\widetilde y$
with
\begin{equation}\label{eq:extracted}
 f^2=D_f,\qquad \tau^2=D_t,\qquad g^2=D_g,
 \qquad \widetilde y^2=D_a.
\end{equation}
The conjugate factor also yields $S\mid A_*$. No adjacent-Pell
descent or nonzero strong ordinate has been used to obtain this
divisibility.

Apply Lemma~\ref{lem:carrier} with $\omega=v-q^2$. It gives
\begin{equation}\label{eq:scaledrecover}
 v=q^2(2v_0+1)\quad(v_0\in\Z),\qquad H\mid K^5.
\end{equation}
Thus $Y=q^3(2v_0+1)$, the exact old cubic scale. In particular
$Y\ne0$. Now if $c=0$, the budget is
\[
 B=k^2\left(q^2(1-4Y^2)-E_0\right)\le0,
\]
contradicting \eqref{eq:budgetpositive}. Hence $c,I\ne0$ and $Q>1$.
This ordering explicitly handles the initially possible branch
$v=c=0$.

\subsection{The auxiliary root is integral}
Put $y=\widetilde y/\delta\in\Q$. Then
\begin{equation}\label{eq:auxiliary}
 QV^2-(Q-1)y^2=1.
\end{equation}
Here $V\in\Z$ and $Q>1$. We recall the elementary descent rather
than silently treating a rational root as integral. In lowest terms
$y=a/b$, $b^2\mid Q-1$. For positive $V>1$ and $y>0$, use
\[
 V'=(2Q-1)V-2(Q-1)y,\qquad
 y'=(2Q-1)y-2QV.
\]
The norm is preserved, $V'$ is integral, and the reduced denominator
of $y'$ remains $b$, since $2Q-1\equiv1\pmod b$. The norm gives
$V<y<QV/(Q-1)$, hence $-V<V'<V$; $V'=0$ is impossible.
Taking absolute values and iterating reaches $|V|=1$, where $|y|=1$.
Thus $b=1$. We have recovered all three ordinary Pell conditions.

\subsection{The inherited rank interface, in the necessary order}
The budget disk implies the common strict signs of $q,F,C,q-F-C$:
\[
 0<\frac{A_0F}{q}<2,\qquad
 0<\frac{B_0C}{q}<2,\qquad
 0<\frac Fq+\frac Cq<1.
\]
The following inherited ordinary interface is the precise input from
the centered signed source\cite{incoming}.

\begin{proposition}[Preprojection rank interface, inherited]
\label{prop:rank}
Assume the authentic six coefficient ports, the packing equations,
$Y=q^3(2v_0+1)$, the common strict signs above, the centered strip
$|c/k-Y|<1/2$, $I\ne0$, $c^2\mid I$, $|I|\ge c^2$, and
\eqref{eq:extracted}--\eqref{eq:auxiliary}. Then, before using any
main projection congruence,
\[
 R=2t+1,\quad t\ge24,\quad
 c=\sigma\psi_{a+2}(R),\quad
 |Y|>\frac{(|X|-3)^t}{2}-1,
 \qquad\sigma=\sgn(c).
\]
Here $\psi_A$ is the Pell polynomial sequence with
$\psi_A(0)=0$, $\psi_A(1)=1$, and
$\psi_A(m+2)=2A\psi_A(m+1)-\psi_A(m)$.
No parity assumption on $I/c^2$ is required.
\end{proposition}

For clarity, its main steps are as follows. Absolute Pell folding first
recovers $|c|=\psi_{|a+2|}(p)$ using
$V\equiv-c\pmod f$, the strong norm, and the auxiliary norm.
The inequality $2|I|<|f|$ prevents wrapping. The first norm supplies
$|k|=2\psi_{|2XY^2+1|}(m)$, and the packing bound places its index
in the required large range. Strong divisibility gives $p\mid m_0$
for the strong index $m_0$. Writing $m_0=pv$ and using
$c^2\mid I$ in the Lucas multiplication formula gives $|c|\mid v$,
so $|c|\mid m_0$. Finally
\[
 V+R=\Delta I^2(c-R)+4ID_fT
\]
implies $V\equiv-R\pmod c$. Signed index reduction and the strict
index bounds identify $p=R$. Pell growth supplies the displayed ratio.
The copied ordinary proofs contain the complete index step-down and
uniform estimates; they are not replaced by finite test cases here.

Let $x=|X|$ and $m=|Y|$. The ratio bound gives
\[
 |H|\ge4(x-1)m-3>x(x-3)^t.
\]
The last inequality follows since
$(x-2)(x-3)^t-4x+1>0$ for $x\ge64$, $t\ge1$.
Set
\begin{equation}\label{eq:projector}
 K_q=X(X+2ac)-1,\qquad
 \Pi=(\sigma X-2^R)(2^R\sigma X+1).
\end{equation}
The exact algebraic identity and inherited signed Pell congruence are
\begin{equation}\label{eq:congruences}
 K=HXc-2K_q,\qquad 2^RK_q\equiv\Pi\pmod H.
\end{equation}
Since $H$ is odd, $H\mid K^5$ implies $H\mid\Pi^5$.
Both factors of $\Pi$ have magnitude below $2^{R+1}x$, so
$|\Pi|<2^{4t+4}x^2$. Lemma~\ref{lem:power5} yields $\Pi=0$.
The second factor is odd and nonzero; therefore
\begin{equation}\label{eq:powerrecovered}
 \sigma X=2^R.
\end{equation}

The inherited centered half-binomial argument now restores positive
orientation and $0<\eta/k<1/4$. One can regenerate fresh odd strong
and auxiliary witnesses for the exact old signed source, keeping the
outer coordinates fixed. Its ordinary native/compiler calibration
then supplies a finite syntactic ZFC refutation. This completes
soundness without assuming the missing main norm during the earlier
root extraction or rank recovery.

\section{Completeness with a fixed strong witness}
\label{sec:completeness}

Start with a finite syntactic ZFC refutation and the canonical outer
data from the same ordinary source theorem. It supplies the cubic
scale $Y=q^3(2v_0+1)$, the correct packing equations, and
\[
 E_0\le q^2/2,\qquad0<\eta/k<1/4,
 \qquad B>q^2k^2/4\ge64,
 \qquad D_t=\tau^2>0.
\]
It also supplies the even Pell base $A=a+2>1$, $R\equiv3\pmod4$,
and $c=\psi_A(R)\equiv R\pmod4$, as proved in Section 7 of the
retained \texttt{centered11-proof.md}. These are the interpolation
hypotheses used below.
For the mixed chart set $v=q^2(2v_0+1)$.

Unconditionally $D_t\equiv1\pmod3$: $\ell=XY^2$ is either $0$
or $2$ modulo $3$. Hence $(3B-4)D_t$ is positive and $2\pmod3$,
so it is not a square. Pell existence supplies $s,g>0$ with
\[
 g^2-(3B-4)D_ts^2=1.
\]
This is exactly $D_g=g^2$; the former squaring of a Pell unit is
unnecessary after deleting the extra factor $4$.

\subsection{Supported-modulus interpolation}
We use the following strengthening proved in the incoming powered
saturation package. We include its mechanism because it is the reason
the factor $s$ can be removed from $I$.

\begin{theorem}[Supported-modulus interpolation, inherited]
\label{thm:interpolation}
For the canonical even Pell base $A=a+2>1$, index $R\equiv3\pmod4$,
and $c=\psi_A(R)>0$ with $c\equiv R\pmod4$, choose one positive
even strong index $m_0$ with
\[
 I=\psi_A(m_0),\quad f=\chi_A(m_0),\quad4sc^2\mid I.
\]
For this same fixed pair $I,f$, every residue class modulo any positive
integer whose prime factors divide $I$ contains a $T$ for which
\[
 V=f^2(c-R+4IT)-c,\qquad
 (\Delta I)^2V^2-((\Delta I)^2-1)y^2=1
\]
has an integer $y$.
\end{theorem}

Here $\chi_A$ is the companion Pell sequence. A finite matrix period
provides the prescribed strong multiple. To see why the modulus can
have larger prime powers than $I$, set
\[
 \rho=\Delta I,\quad \lambda=4m_0f,
 \quad d=\operatorname{lcm}(\lambda,4I).
\]
The squared-modulus period is $\lambda$, not merely $4m_0$.
For $p\mid I$, the Lucas valuation inequality
$v_p(m_0)\le v_p(I)$ and $\gcd(f,I)=1$ imply
\[
 v_p(d)=v_p(4I).
\]
Start from one admissible odd auxiliary index $\nu_0\equiv3\pmod4$
and vary it through $\nu_0+dh$ for $h\ge0$. With
$V_\nu=\chi_\rho(\nu)/\rho$, put
\[
 T(h)=\frac{V_{\nu_0+dh}+c-f^2(c-R)}{4If^2}.
\]
These are integers. Pell addition and elementary lifting give the
exact isometry
\begin{equation}\label{eq:isometry}
 v_p\bigl(T(h_1)-T(h_2)\bigr)=v_p(h_1-h_2)
 \qquad(p\mid I,\ h_1,h_2\ge0,\ h_1\ne h_2).
\end{equation}
For example, for $d_0$ divisible by $4$ and $p\mid\rho$,
\[
 v_p\left(\psi_\rho(d_0)/\rho\right)=v_p(d_0),
\]
while the other summand in the Pell addition formula has strictly
higher valuation. The formulas follow by expanding powers of the
fourth unit, whose ordinate is $4\rho(2\rho^2-1)$; the doubling
formula handles $p=2$. Dividing by $4If^2$ proves
\eqref{eq:isometry}. It makes the map bijective modulo every $p^e$
with $p\mid I$, and CRT proves the theorem. The full proof,
including the two-adic valuation and the base auxiliary index, is
retained as \texttt{inherited/norm/powered\_saturation\_proof.tex}.

\subsection{Finish the witness construction}
Choose the strong pair once as in Theorem~\ref{thm:interpolation} and
set
\[
 r=I/c^2>0.
\]
Every prime factor of $sr$ divides $I$. In the mixed chart the
canonical projection and scale give
\[
 \kappa=\frac K H+\frac{v-q^2}{M}\in\Z,
 \qquad A_*=(MH)^5\pi(T+\kappa).
\]
Intersectivity supplies a residue $a_0$ modulo $sr$ for which
$\pi(a_0+\kappa)\equiv0\pmod{sr}$. Only now request the auxiliary
coordinate $T\equiv a_0\pmod{sr}$ from the already fixed strong
pair. It follows that
\[
 \frac{A_*}{S}=\frac{H^4\pi(T+\kappa)}{sr}\in\Z.
\]
Choose
\begin{equation}\label{eq:ncomplete}
 n=A_*/S+f+3\delta y+9\tau g-3QV.
\end{equation}
One conjugate factor is zero, proving completeness. The cubic chart
uses $\kappa=K/H$ and the identical argument. The direct companions
omit the last term in \eqref{eq:ncomplete}.

The witness order is important: outer native data, then $s,g$, then
one fixed strong pair $I,f,r$, then an intersective residue, then
$T,y$, and finally $n$. No modulus is repaired by changing $I$
after a residue has been chosen.

\begin{theorem}[Eleven-integer consistency equations]\label{thm:main}
Each of the four supplied eleven-integer polynomials has an integer
zero if and only if the pinned ordinary source has a finite ZFC
refutation. Composed with that source's ordinary calibration, ZFC
proves \eqref{eq:goal} for each polynomial. Their exact degrees and
arithmetic lengths are those in Table~\ref{tab:norm}.
\end{theorem}
\begin{proof}
The two existence directions were established in
Sections~\ref{sec:soundness} and~\ref{sec:completeness}. All new
arguments are elementary arithmetic, Pell identities, finite group
periods, and CRT, formalizable in ZFC. The next section proves exact
degree, and the finite circuit audit recounts arithmetic length.
Negating the finite-refutation existence statement yields the claimed
equivalence. This does not establish either side of that equivalence.
\end{proof}

\section{Exact degree, cancellation, and arithmetic length}
\label{sec:degree}

The degree reductions are polynomial identities, not degree estimates
from evaluations at random points. Table~\ref{tab:degrees} records the
relevant exact degrees; $B_*=Sn-A_*$ and $C_0=3SQV$.

\begin{table}[ht]
\centering
\begin{tabular}{lrrrrrrrr}
\toprule
Chart & $S$ & $D_f$ & $Q$ & $V$ & $D_a$ & $P$ & $A_*$ & $C_0$\\
\midrule
Cubic & 8 & 18 & 30 & 23 & 106 & 94 & 35 & 61\\
Mixed & 16 & 14 & 22 & 19 & 82 & 62 & 35 & 57\\
\bottomrule
\end{tabular}
\caption{The mixed budget has degree $16$, not $14$: its
$-4q^2\eta^2$ term dominates. This affects $P$ and the final degree.}
\label{tab:degrees}
\end{table}

The exact identity
\begin{equation}\label{eq:cheapcenter}
 b-C_0^2=9S^2\bigl(-Q(V^2+1)+1\bigr)
\end{equation}
shows that $\omega=Z^2-b$ has degree at most $96$ on the cubic
chart and at most $92$ on the mixed chart. The full three-root norm
can be rewritten without changing the polynomial as
\begin{equation}\label{eq:degreecut}
 \left[(\omega-u-e)^2-4\{ue+b(u+e)\}\right]^2
                  -64ubeZ^2.
\end{equation}
For the mixed chart, $\deg(u,b,e)=(46,114,94)$. The unique highest
part of the inner bracket is $-4b_0e_0$, of degree $208$; the outer
correction has degree at most $368$. Thus the exact leading form is
$16b_0^2e_0^2$ and the degree is $416$. For the cubic chart,
$\deg(u,b,e)=(34,122,110)$, giving degree $464$; its outer
correction has degree at most $388$.

For an explicit nonvanishing certificate at the actual fixed tuple,
the mixed polynomial has
\begin{equation}\label{eq:leadmixed}
 [j^{216}w^{52}v^{64}h^{12}z^4s^{12}r^{24}c^{32}]P_{416}
   =2^{384}3^{282}\beta^{216}>0.
\end{equation}
The cubic centered polynomial has
\begin{equation}\label{eq:leadcubic}
 [j^{264}w^{52}v^{64}h^{12}z^4s^{12}r^{24}c^{32}]P_{464}
   =2^{456}3^{330}\beta^{264}>0.
\end{equation}
The exponent sums are respectively $416$ and $464$.

One way to verify the mixed certificate is to put
\[
 q_0=6\beta j,\quad X_0=12q_0w,\quad Y_0=q_0v,
 \quad R_0=-q_0^3(12wC-q_0z),\quad k_0=R_0+hX_0Y_0.
\]
The highest homogeneous form is
\[
 2^{64}3^{14}s^{12}r^{24}c^{32}
 q_0^{84}X_0^{40}Y_0^{52}R_0^4k_0^{12}.
\]
Selecting $h^{12}$ in the last factor and $z^4$ in $R_0^4$ gives
\eqref{eq:leadmixed} uniquely. The cubic calculation uses
$Y_0=2vq_0^3$, $k_0=hX_0Y_0$, and the corresponding leading form
with coefficient $2^{24}3^{14}$ and $q_0^4$. Thus no possible
cancellation involving one of the other five fixed ports can affect
the nonzero witness.

Without the center, $Z$ has degree $35$ and $b$ is the unique
dominant radicand. The leading term is $b_0^4$, giving degrees
$456$ and $488$. The finite auditor checks the complete norm
identity, these cancellation cuts, the leading certificates, all
declared variables, and the actual graph counts.

The principal graph has $24$ additions, $22$ subtractions, and $83$
multiplications, totaling $129$. The centered cubic graph has
$24+21+78=123$. The simpler center $3QV$ is sufficient for the exact
degree reduction. Subtracting another multiple of $V$ would refine an
approximation that is not responsible for the highest remaining degree.

\section{A direct quartic in 478 integer variables}
\label{sec:direct478}

The preceding compact equations optimize the number of integer unknowns.
There is also a concrete improvement at degree four using the separate
direct C222 proof-certificate source.  Let $R_0,\ldots,R_{304}$ be the
literal integral quadratic polynomials in
\texttt{direct/equation.json}, in its specified coordinate order, and set
\begin{equation}\label{eq:direct478}
 \boxed{P_{478}(z_0,\ldots,z_{477})=\sum_{j=0}^{304}R_j(z)^2.}
\end{equation}
The companion \texttt{expanded\_polynomial.json} contains every collected
coefficient.  A term $[c,[i_1,\ldots,i_d]]$ denotes
$cz_{i_1}\cdots z_{i_d}$, with an empty index list denoting the constant
monomial.  Every coefficient is a literal integer, and every variable
ranges independently over $\Z$.

\begin{theorem}\label{thm:direct478}
The polynomial \eqref{eq:direct478} has exact degree four, 478 occurring
variables, 5,484 collected monomials and coefficient height 511,326.
The ordinary proofs below, composed with the explicitly retained direct
C222/refutation calibration, establish
\[
 \mathrm{ZFC}\vdash
 \bigl[\neg\exists z\in\Z^{478}\ P_{478}(z)=0\bigr]
 \ \longleftrightarrow\ \operatorname{Con}(\mathrm{ZFC}).
\]
\end{theorem}

The matched predecessor is the derived-label square-shell endpoint in the
incoming powered-saturation/factorial-carrier archive at the comparison
snapshot.  Its metrics, and the new metrics, are:
\begin{center}
\begin{tabular}{lrrrr}
 & Variables & Quadratic rows & Monomials & Height\\
\hline
Matched predecessor & 517 & 310 & 6,774 & 3,904,267\\
New direct equation & 478 & 305 & 5,484 & 511,326
\end{tabular}
\end{center}
This improves all four displayed metrics for that direct branch.  Other
repository constructions have fewer variables with different degree or
coefficient costs.  No global minimum is asserted.

\subsection{The dependency order behind the saving}

The source supplies a proposed factorial $F$, a horizon $t=t_0+16$,
a proposed power $V$, and exact polynomial relations
\begin{align}
 &D=r+s+1, & DF&=V+s,\\
 &q=D+o_0+1, &q&=B+o_1+1,\\
 &B i=C+F^2+1, &&C\ge1.
\end{align}
Here $D$ and $B$ are proposed binomial outputs.  Their eventual meanings
are not assumed in these identities.  The power interface requests
\[
 V=(Ft^2+1)^t,\qquad N=B^K,\qquad q=N(N+V+1)^K,
 \qquad K=2(2+t+x+y+F),
\]
as well as the powers used for digit and range extraction.
Before power recovery, retain $t_0,x\ge0$, the endpoint and encoded-input
signs, and $o_0\ge0$.  Strengthen the existing factorial guard to $F\ge1$.
It costs the same number of witnesses and holds for every genuine factorial.

\begin{lemma}[Factorial signs after one power]\label{lem:early-fact-signs}
Suppose $F\ge1$, $V>F$, $D=r+s+1$, $DF=V+s$ and $rs\ge0$.
Then $r,s\ge0$ and $D\ge1$.
\end{lemma}
\begin{proof}
Eliminate $D$ to obtain $V=Fr+(F-1)s+F$.  If either member of $(r,s)$
were negative, product nonnegativity would make both nonpositive.
The displayed expression would then give $V\le F$, a contradiction.
This includes the cases where one member is zero.
\end{proof}

The base $Ft^2+1$ is already at least two.  The signed power certificate,
whose frame is justified below independently of these signs, therefore
recovers $V=(Ft^2+1)^t>F$ first.  Lemma~\ref{lem:early-fact-signs} proves
$D\ge1$, and hence $q=D+o_0+1\ge2$.  Next require only $Bo_1\ge0$.
The sum $B+o_1=q-1>0$ forces $B,o_1\ge0$; the reciprocal index identity
excludes $B=0$, so $B\ge1$ before its power is interpreted.

The two calls at index $K$ now orient their shared Pell pair.  The first
either gives $N=B^K$ or has the exceptional inverse case $B=1,N=-1$.
In both cases its second base $N+V+1$ is at least two.  That second call
fixes the orientation and excludes the inverse case.  All remaining seed
pairs have supporting calls at $q$ or $q+1$, both at least two.  The
unchanged addition formula recovers the single derived $t+1$ pair.
Thus all nineteen requested powers have their intended values.

Only now recover the other delayed signs.  Both binomial remainder pairs
and all ten range remainder pairs have exact sum equations
\[
 r_j+s_j=q^{\ell_j}-1>0,
\]
where the already positive exponent is $t,t+1,x+1$, or $y+1$.
For each such pair, $r_js_j\ge0$ implies both coordinates are nonnegative.
This replaces fourteen pairs of scalar predicates by fourteen product
predicates.  No binomial, factorial or CRT interpretation is used to
restore a sign that its own interpretation needs.

For unrestricted integer witnesses, use
\[
 Q(a,b,c)=a^2+b^2+c(c+1).
\]
It represents exactly the nonnegative integers.  Nonnegativity is
immediate, and surjectivity follows from the three-square theorem applied
to $4n+1$: exactly one of the three squares is odd.
Consequently $rs=Q(a,b,c)$ is one quadratic row expressing $rs\ge0$.
The ordinary positive-sum/product observation is elementary and already
used in earlier repository work; the new saving is its dependency-ordered
application to this stronger source, enabled by the earlier factorial
bootstrap.

\subsection{One guard for orientation and a nonzero shared ordinate}

Let $A$ be the main Pell base, and write the shared auxiliary ordinate as
\[
 V_* = \left(\prod_iY_i^2\right)(v+1).
\]
The retained seed slacks give $Y_i\ge k_i\ge1$ before any power is
interpreted.  Separate guards on the frame orientation and on $v$ can be
replaced by two quadratic rows involving four fresh coordinates:
\begin{equation}\label{eq:joint-anchor}
 u=(v+1)^2,\qquad
 (A-1)u-1=\alpha^2+\beta^2+\gamma(\gamma+1).
\end{equation}
Since the right side is nonnegative, $(A-1)u\ge1$.  Thus $A\ge2$,
$u\ge1$ and $v\ne-1$, so $V_*\ne0$.  The sign of $V_*$ is immaterial.
Conversely, canonical positive Pell witnesses have $A\ge2,v\ge0$;
the left side in the second equation is nonnegative and admits the required
three-square witnesses.  This joins two useful conclusions in one predicate
while keeping the final equation quartic.

The unchanged auxiliary norm gives
\[
 U^2-(A^2-1)V_*^2=1,\qquad |U|\ge A.
\]
The comparison base $b$ satisfies
\[
 b=1+4V_*(q_4+1),\qquad b\equiv A\pmod U.
\]
We permit $q_4$ to have either sign.  The first equation gives
$b\equiv1\pmod4$; the second excludes $b=1$ because
$0<A-1<|U|$.  Therefore $|b|\ge3$.

The signed comparison-base index theorem is inherited from the incoming
quadratic-frames package, and its complete proofs are retained under
\texttt{direct/source/}.  For completeness, the mechanism is as follows.
Classify absolute Pell solutions at $A$ and signed solutions at $b$;
polynomial dependence in the base transfers their congruences.
Squared divisibility $Y_i^2\mid V_*$ gives $Y_i\mid h$ for the common
auxiliary index $h$.  The abscissa residue identity gives the comparison
index congruent to either sign of the main index modulo $2h$.
The base-one ordinate congruence gives either sign of $k_i$ modulo $4Y_i$.
Hence the main index is congruent to $\pm k_i$ modulo $2Y_i$.
Both lie in $[1,Y_i]$, so they agree; the possible common endpoint also
gives equality.  Thus
\[
 |X_i|=\chi_A(k_i),\qquad Y_i=\psi_A(k_i).
\]
This argument allows both signs of $V_*,U,b$ and the comparison coordinates.
It never requires the removed private signs to be pointwise positive.

\subsection{A frame that uses proved signs efficiently}

The original power frame squared every atom.  The new frame uses
\begin{equation}\label{eq:typed-frame}
 W=2+\sum_{a\text{ already proved }\ge0}a
       +\sum_{a\text{ of undetermined sign}}a^2.
\end{equation}
There are 41 distinct affine atoms: 18 enter linearly and 23 quadratically.
Their signs are certified by 163 ordered polynomial deductions from 31
retained scalar predicates.  The frame-definition row and every power
interpretation are excluded from those deductions.
For every integer $a$, $a^2\ge|a|$; hence $W\ge2$ and $W>|a|$ for every
base, exponent and proposed output in the finite power interface.

Writing $Z=Z_{\rm minus1}+1$, the unchanged frame norm is
\[
 A^2-W^3(W+2)Z^2=1.
\]
Use $\chi_a(j)+\psi_a(j)\sqrt{a^2-1}=(a+\sqrt{a^2-1})^j$.
Equation~\eqref{eq:joint-anchor} gives $A\ge2$, so $Z\ne0$ even though
its sign is unrestricted.  Pell classification at base $W+1$ and the
recurrence $\psi_{W+1}(h)\equiv h\pmod W$ yield
$A=\chi_{W+1}(h)$ with $h\ge W$.  Thus
\[
 A\ge(W+1)^W,\qquad
 T_n=2An-n^2-1>W^W+W \quad(1\le n<W).
\]
For example $T_n\ge2A-2$, and the binomial expansion gives
$(W+1)^W\ge2W^W$.  These bounds hold before any proposed output is
interpreted or given a sign.

At an indexed pair, the power row
\[
 X+(n-A)Y=m+T_nw
\]
gives $m\equiv n^k\pmod{T_n}$ for positive $X$, and
$mn^k\equiv-1\pmod{T_n}$ for negative $X$.  The frame bounds give
$|m-n^k|<T_n$ and $|mn^k+1|<T_n$.  A base $n\ge2$ excludes the
inverse case and gives $m=n^k$ exactly, with no prior sign assumption
on $m$.  This licenses the initial factorial call and the dependency order
used above.

\subsection{Both directions and the finite equation}

Given a zero of the new system, the preceding arguments recover all powers
and all semantic front signs.  The inherited self-bounding factorial,
exact digit, reciprocal prime-power CRT, derived-label bound and
pairwise-coprime Horn-flow proofs then recover a finite C222 refutation.
Their full ordinary proof texts are included under \texttt{direct/source/}.
The endpoint composition is
\texttt{KEndpoint.c222\_iff\_zf\_refutation} followed by
\texttt{ConstructibleChoice.zf\_refutation\_iff\_zfc\_refutation}.

Conversely, start with the inherited arithmetic front supplied by any
finite refutation.  It has $F=t!\ge1$ and nonnegative entries in every
fused pair.  Construct its new $W$ by \eqref{eq:typed-frame}, choose the
positive Pell solution at index $W$ for base $W+1$, and choose the main
pairs at their required indices.  The order of a unit in a finite
quadratic ring supplies a common positive auxiliary iterate divisible by
the entire product $\prod_iY_i^2$.  CRT supplies a positive comparison
base with the required two congruences.  Complete all comparison pairs
and power quotients.  These canonical witnesses have $v\ge0$, so
\eqref{eq:joint-anchor} and every product predicate admit their ordinary
three-square witnesses.  No consistency assumption is used in this map.

These are equivalences of existential zero assertions.  New private
comparison coordinates need not solve the old positive-domain auxiliary
rows; fresh old witnesses can be reconstructed after recovering the front.

The complete new source has 690 integer coordinates and 515 quadratic
rows.  The projection certificate records 210 integral monic restorations;
two remaining unused coordinates are discarded with restoration zero.
Every substitution keeps every residual quadratic.  This leaves 478
coordinates and 305 rows.  The final explicitly recorded integral
translations are bijections.  Expanding the sum of squares gives all 5,484
nonzero terms of \eqref{eq:direct478} and proves the stated exact metrics.

The portable builder regenerates all ten normative data files byte for
byte.  One independent checker reconstructs every guard and frame row;
a second checks all source hashes, transformations, projection inverses,
translations and collected quartic coefficients.  Both pass under normal
and optimized Python.  Their finite boundary controls support transcription
and degeneration checks; they neither decide solvability nor replace the
unbounded mathematical proof.  This branch uses the direct calculus,
without a K215 execution encoding or Jones universal-polynomial parameter.


\section{Fixed numeral definitions}
\label{sec:constants}

The eleven-variable forms retain one complete fixed source34 tuple.
It is deliberately separate from the direct quartic's small literal
coefficient list. The byte files are actual numerals: a decompressed
little-endian payload $(b_0,\ldots,b_{m-1})$ denotes exactly
\[
 \sum_{i=0}^{m-1} b_i256^i.
\]
Hashes authenticate those bytes and do not stand in for unspecified
integers. The six minimal payloads total $615{,}883$ raw bytes.

\begin{center}
\begin{tabular}{lrr}
\toprule
Port & Bit length & Raw bytes\\
\midrule
$\beta$ & 1064921 & 133116\\
$K_0$ & 1065043 & 133131\\
$M_C$ & 1064921 & 133116\\
$M_F$ & 1064922 & 133116\\
$A_0$ & 1114 & 140\\
$B_0$ & 666105 & 83264\\
\bottomrule
\end{tabular}
\end{center}

For the finite recipe, put $\rho=2^{121}$ and
$\mathcal B=\rho^{8801}=2^{1064921}$. The retained file
\texttt{wang943\_tiles.json} gives the ordered $943$ tiles in
west--north--east--south order, reconstructed from the included
literal $430$ transitions. Start is selector $0$ and forbidden is
selector $1$. Let $x_{it}$ be the four labels of tile $i$, and put
\begin{align*}
 W_j&=2^{10+2j},\qquad (o_0,o_1,o_2,o_3)=(1020,1000,980,960),\\
 \mathcal E&=\{0,\ldots,942\}\cup
 \{o_t+2j:0\le t<4,\ 0\le j<10\}\\
 &\hspace{12mm}\cup\{1063,1064,1095,1096,2195,3296\},\\
 c_i&=1+\sum_{t=0}^3\sum_{j=0}^9
       \operatorname{bit}_j(x_{it}+512)W_{10t+j}
       +\sum_{j=40}^{43}W_j\quad(0\le i\le942).
\end{align*}
At the copy exponent $o_t+2j$ put $c_{o_t+2j}=W_{10t+j}$; at
$1095,1096,2195,3296$ put $W_{40},W_{41},W_{42},W_{43}$ respectively.
The dummy coefficients at $1063,1064$ are zero. Then
\begin{align*}
 \mu&=1022+\sum_{j=0}^{43}W_j,\\
 Z_p&=(2^{119}-1)\sum_{e=8792}^{8798}\rho^e
                  +(2^{27}-1)\rho^{8799},\\
 D_C&=\rho^{40}+\rho^{41}+\rho^{1100}+\rho^{2200}
       +\sum_{e\in\mathcal E}c_e(\rho^{4393-e}+\rho^{5495-e}),\\
 M_C&=\mathcal B-1-\sum_{e\in\mathcal E,\ e\ne1}\rho^e
                         -2\rho^{1063},\\
 M_{F,n}&=4+\mu(\rho^{4393}+\rho^{5495})
       +\sum_{j=0}^9(\rho^{1000+2j}+\rho^{1021+2j})\\
       &\hspace{12mm}+\rho^{2195}+\rho^{3296}+Z_p.
\end{align*}
The ports are precisely
\[
 \beta=\mathcal B-1,\quad K_0=D_C+\mathcal B\rho,
 \quad M_C,\quad M_F=M_{F,n}+\mathcal B-1,
 \quad A_0=2^{1113},\quad B_0=2^{666104}.
\]
The authentic tuple has an ordinary compiler adapter and canonical
rank controller, retained in the source material. Matching only
these residues or substituting a different tuple's mask would not
justify the consistency interpretation.

\section{Artifact design, reproduction, and verification boundaries}
\label{sec:verification}

The archive contains the full manuscript source and PDF, the four
literal eleven-variable graphs, the direct residual and expanded
polynomial files, deterministic builders, independent checkers,
fixed numeral payloads, and scoped inherited proof texts. The
top-level \texttt{README.md} gives exact commands and the final
file map. No network lookup is required to run the new finite checks.

\paragraph{Small-variable graph format.}
Nodes are topologically ordered integer, coefficient, or variable
leaves, followed by binary \texttt{add}, \texttt{sub}, or
\texttt{mul} nodes referring to earlier node indices. The output
index identifies one polynomial. In particular, register names do
not add variables. A reader can reconstruct the expression without
executing the builder or trusting its operation count.

\paragraph{Direct sparse format.}
Each residual is an explicit quadratic integer polynomial. The
full quartic is its collected sum of squares, stored separately and
re-expanded independently. Every coefficient and variable occurrence
is supplied; the final polynomial is not a request to call a proof
predicate, a binomial oracle, or a universal interpreter. Arithmetic
graphs used during its derivation have already been expanded into
ordinary polynomial rows.

\paragraph{What the checks establish.}
Exact algebra checks transcription, the $35$-term norm identity,
degree cancellations, literal coefficients, and the direct sum of
squares. Graph checks establish variable closure, liveness and
counts. Source and coefficient hashes establish byte identity.
Reconstruction checks reproduce the supplied new artifacts in a
separate output directory. Selected deliberate corruptions must be
rejected. Optimized-Python runs check that acceptance does not
depend on disabled \texttt{assert} statements. The machine-readable
receipts state their actual scopes.

These are useful finite certificates. They do not by themselves
establish universal Pell classification, supported-modulus
interpolation, the native compiler, the C222 calibration, or
consistency. Those are mathematical dependencies and ordinary
proofs, identified separately. No search for a ZFC contradiction,
nor a bounded failure to find an integer solution, is used as
evidence of \eqref{eq:goal}.

\section{Further research questions}
\label{sec:further}

\begin{enumerate}
\item \textbf{Smaller mixed carriers.} The current mixed chart needs
      $M^5\mid A_*$ to recover $M\mid v-q^2$ by a least-valuation
      argument. The $H$ part only needs first-power divisibility.
      A carrier tailored to the supported strong modulus might reduce
      the $M^5$ cost, but its completeness must work at every prime
      dividing the chosen strong ordinate, not just at generic primes.
\item \textbf{An eleven-integer equation below degree $416$.}
      The leading obstruction is the mixed term $16b_0^2e_0^2$.
      Refining the auxiliary center alone does not remove it.
      Reducing the budget/guard degree, or proving a different fusion
      between the guard and a Pell condition, is a more consequential
      target than another correction to the same center.
\item \textbf{Ten integers without the general reduction.}
      The guard parameter $s$ and strong multiplier $r$ are still
      independent. Identifying them naively imposes simultaneous Pell
      conditions with a common ordinate; ordinary Pell existence
      does not prove that those simultaneous conditions have a
      solution. A valid elimination must solve that completeness
      problem explicitly.
\item \textbf{Further direct sign sharing.}
      The $478$-variable quartic uses exact source identities to
      recover individual signs from shared products. Further savings
      require a proved positive sum or a similarly noncircular
      certificate. The number of removed guards is not the final
      variable saving, since a product can retain coordinates that
      affine projection previously removed.
\item \textbf{Lower coefficients in the direct route.}
      The new frame lowers height substantially on its matched
      branch, but more variables can still buy smaller coefficients.
      An explicit Pareto search should measure collected terms and
      height after projection, not just gate counts before projection.
\item \textbf{Formalize one complete bridge.}
      The most valuable formal target is the ordinary equivalence
      between one fully emitted direct polynomial and finite C222
      refutability, connected to the already identified endpoints.
      Local norm identities alone would not close that logical gap.
\end{enumerate}

\section*{Attribution and intended integration}
The inherited proof and code notices are retained. New prose, circuits,
and checks were prepared with AI assistance for incorporation into
Vladimir Reshetnikov's ProveIt repository. The work builds on the
specific incoming and repository results identified above; classical
Pell arithmetic, CRT, norm identities, and intersective polynomials
are credited as inherited tools. No repository branch was modified
or submission published by producing this archive.

{\raggedright
\begin{thebibliography}{9}
\bibitem{repo}
Vladimir Reshetnikov, \emph{ProveIt}, source snapshot
\texttt{45b399ec51df}\allowbreak\texttt{becb891dd51d8a9e4b6b7c361ac6},
24 October 2026.
\url{https://github.com/VladimirReshetnikov/ProveIt}.
The relevant directories are \texttt{Computability/DiophantineConsistency},
\texttt{Computability/HilbertTenthProblem}, and \texttt{docs/incoming}.

\bibitem{incoming}
ProveIt incoming research packages, \emph{Coprime Fusion and Low-Degree
Diophantine Projections}; \emph{Powered Saturation and Factorial
Carriers, Round 14}; \emph{Weaker Congruences and Even Pell Indices};
and \emph{Quadratic Frames}, 24 October 2026.
Exact archive names, source bytes, and dependency roles are recorded
in this release's provenance ledgers.
\url{https://github.com/VladimirReshetnikov/ProveIt/tree/45b399ec51dfbecb891dd51d8a9e4b6b7c361ac6/docs/incoming}.

\bibitem{hyde2014}
A. M. Hyde, P. D. Lee, and B. K. Spearman,
\emph{Polynomials $(x^3-n)(x^2+3)$ solvable modulo any integer},
American Mathematical Monthly \textbf{121}(4) (2014), 355--358.
\url{https://doi.org/10.4169/amer.math.monthly.121.04.355}.

\bibitem{sun2026}
Zhi-Wei Sun, \emph{Ten unknowns for Hilbert's tenth problem over
the integers}, arXiv:2609.01594v2, 2 September 2026.
\url{https://arxiv.org/abs/2609.01594}.
The present eleven-variable construction does not instantiate that
general ten-variable reduction.
\end{thebibliography}
}
\end{document}
